\section{Conclusion}
\label{sec:conclusion}

This review has provided a systematic examination of signal processing methods for lubrication condition monitoring of ball bearings, the first review in the LCM literature to organise its findings by signal processing method family rather than by sensing modality. A total of 48 primary sources were reviewed against inclusion criteria requiring online or in-line data acquisition, lubrication-specific monitoring objectives, and sufficient methodological detail to assess the signal processing approach.

On the sensing side, the review establishes a two-axis classification -- active/passive and invasive/non-invasive -- and shows that passive non-invasive methods (acoustic emission, vibration, passive ultrasound, microphone arrays, thermal monitoring) carry genuine lubrication-relevant information through the physical mechanisms summarised in Table~\ref{tab:signal_origins}. All of these signals increase in amplitude as the lubrication film thins and asperity contacts increase -- a property that makes them inherently usable as lubrication condition indicators -- but their amplitude is confounded by operating speed, load, and temperature independently of lubrication state. Active and invasive methods (ultrasonic reflectometry, electrical capacitance/impedance, optical, magnetic) provide more direct film thickness information via physics-based models, at the cost of installation complexity and, in some cases, bearing modification.

On the signal processing side, the master cross-reference matrix (Table~\ref{tab:sp_crossref}) reveals a wide diversity of methods applied to LCM, spanning from computationally trivial time-domain statistics (RMS, pulse count rate) to sophisticated methods (waveform clustering, graph neural networks, physics-informed $\kappa$ regression). Time-domain statistical features dominate practical and industrial applications due to their low computational cost and robustness. Frequency-domain methods -- particularly envelope analysis, spectral kurtosis, and fractional-octave filter banks -- add mechanistic discrimination that time-domain methods cannot provide, enabling discrimination between starvation, contamination, and other lubrication fault types. Physics-based model-driven methods (spring-layer acoustic model, electrical circuit models, $\kappa$/$\Lambda$ regression targets) bridge signal processing and lubrication physics, enabling quantitative film thickness estimation from sensor signals. Machine learning and data fusion methods show diagnostic promise but are at the research stage.

The comparative analysis in Section~\ref{sec:comparative} reaches five principal conclusions: passive non-invasive methods are most sensitive to starvation, least sensitive to contamination, and confounded by operating condition variation; active methods provide quantitatively superior film thickness information; thermal monitoring responds more slowly than acoustic methods; no SP method has been validated under realistic variable-speed, variable-load industrial conditions; and the technology maturity gap between research and industry is large, with all sophisticated SP approaches remaining at the research or prototype stage.

Six signal-processing-centric research gaps are identified. The most consequential are the absence of standardised benchmark datasets that would enable rigorous cross-study comparison, the unsolved problem of speed- and load-invariant feature extraction for passive methods, and the underexploration of grease lubrication monitoring despite grease being the dominant lubricant in industrial rolling element bearings. Additional gaps include the disconnect between condition assessment outputs and actionable relubrication decisions, the lack of multi-sensor fusion validation in realistic industrial environments, and the untapped potential of self-supervised and foundation model approaches for learning informative representations from the large archives of unlabelled bearing monitoring data that exist in industrial settings.

The overarching conclusion is that the signal processing layer -- not the sensing hardware -- is the primary barrier limiting LCM from laboratory demonstration to industrial deployment. Advanced sensors capable of detecting lubrication-relevant signals exist and are demonstrated. What is lacking are SP methods that are simultaneously physically interpretable, robust to operating condition variation, validated under industrial conditions, and integrated with downstream maintenance decision systems.
